\subsection{THE EXCHANGE CONDITION}

The ``exchange condition'' is the following assertion about (W, S):

\begin{enumerate}
\def\labelenumi{(\Alph{enumi})}
\setcounter{enumi}{4}
\tightlist
\item
  Let \(w \in W\) and \(s \in S\) be such that \(l(sw) \leqslant l(w)\) . For any reduced decomposition \((s_1, \ldots, s_q)\) of \(w\) , there exists an integer \(j\) such that \(1 \leqslant j \leqslant q\) and
\end{enumerate}

\[
s s _ {1} \dots s _ {j - 1} = s _ {1} \dots s _ {j - 1} s _ {j}.\tag{16}
\]

We assume in this number that \((W,S)\) satisfies (E). By Lemma 3, this is so if \((W,S)\) is a Coxeter system. The results of this number thus apply to Coxeter systems.

PROPOSITION 4. Let \(s \in S\) , \(w \in W\) and \(s = (s_1, \ldots, s_q)\) be a reduced decomposition of \(w\) . Then one of the following must hold:

\begin{enumerate}
\def\labelenumi{\alph{enumi})}
\item
  \(l(sw) = l(w) + 1\) and \((s, s_1, \ldots, s_q)\) is a reduced decomposition of sw.
\item
  \(l(sw) = l(w) - 1\) and there exists an integer \(j\) such that \(1 \leqslant j \leqslant q\) , \((s_1, \ldots, s_{j-1}, s_{j+1}, \ldots, s_q)\) is a reduced decomposition of \(sw\) and \((s, s_1, \ldots, s_{j-1}, s_{j+1}, \ldots, s_q)\) is a reduced decomposition of \(w\) .
\end{enumerate}

Put \(w' = sw\) . By formula (3) of no. 1, we have

\[
\left| l (w) - l \left(w ^ {\prime}\right) \right| \leqslant l (s) = 1.
\]

We distinguish two cases:

\begin{enumerate}
\def\labelenumi{\alph{enumi})}
\tightlist
\item
  \(l(w') > l(w)\) . Then \(l(w') = q + 1\) and \(w' = ss_1 \ldots s_q\) , so
\end{enumerate}

\[
\left(s, s _ {1}, \ldots , s _ {q}\right)
\]

is a reduced decomposition of \(w'\) .

\begin{enumerate}
\def\labelenumi{\alph{enumi})}
\setcounter{enumi}{1}
\tightlist
\item
  \(l(w') \leqslant l(w)\) . By (E), there exists an integer \(j\) such that \(1 \leqslant j \leqslant q\) and (16) is satisfied. Then \(w = ss_1 \ldots s_{j-1}s_{j+1} \ldots s_q\) and hence
\end{enumerate}

\[
w ^ {\prime} = s _ {1} \dots s _ {j - 1} s _ {j + 1} \dots s _ {q}.
\]

Since \(q - 1 \leqslant l(w') \leqslant q\) , it follows that \(l(w') = q - 1\) and that \((s_1, \ldots, s_{j-1}, s_{j+1}, \ldots, s_q)\) is a reduced decomposition of \(w'\) .

Lemma 4. Let \(w \in W\) have length \(q \geqslant 1\) , let \(D\) be the set of reduced decompositions of \(w\) , and let \(F\) be a map from \(D\) to a set \(E\) . Assume that \(F(s) = F(s')\) if the elements \(s = (s_1, \ldots, s_q)\) ; \(s' = (s_1', \ldots, s_q')\) of \(D\) satisfy one of the following hypotheses:

\[
a) s _ {1} = s _ {1} ^ {\prime} \text {   or   } s _ {q} = s _ {q} ^ {\prime}.
\]

\begin{enumerate}
\def\labelenumi{\alph{enumi})}
\setcounter{enumi}{1}
\tightlist
\item
  There exist \(s\) and \(s'\) in \(S\) such that \(s_j = s'_k = s\) and \(s_k = s'_j = s'\) for \(j\) odd and \(k\) even.
\end{enumerate}

Then F is constant.

\begin{enumerate}
\def\labelenumi{\Alph{enumi})}
\item
  Let \(\mathbf{s}, \mathbf{s}' \in D\) and put \(\mathbf{t} = (s_1', s_1, \ldots, s_{q-1})\) . We are going to show that if \(\mathrm{F}(\mathbf{s}) \neq \mathrm{F}(\mathbf{s}')\) then \(\mathbf{t} \in D\) and \(\mathrm{F}(\mathbf{t}) \neq \mathrm{F}(\mathbf{s})\) . Indeed, \(w = s_1' \ldots s_q'\) and hence \(s_1' w = s_2' \ldots s_q'\) is of length \(< q\) . By Prop. 4, there exists an integer \(j\) such that \(1 \leqslant j \leqslant q\) and the sequence \(\mathbf{u} = (s_1', s_1, \ldots, s_{j-1}, s_{j+1}, \ldots, s_q)\) belongs to D. We have \(\mathrm{F}(\mathbf{u}) = \mathrm{F}(\mathbf{s}')\) by condition \(a\) ; if \(j \neq q\) we would have \(\mathrm{F}(\mathbf{s}) = \mathrm{F}(\mathbf{u})\) for the same reason, and hence \(\mathrm{F}(\mathbf{s}) = \mathrm{F}(\mathbf{s}')\) contrary to our hypothesis. Thus \(j = q\) and hence \(\mathbf{t} = \mathbf{u} \in D\) and \(\mathrm{F}(\mathbf{t}) = \mathrm{F}(\mathbf{s}') \neq \mathrm{F}(\mathbf{s})\) .
\item
  Let \(\mathbf{s}\) and \(\mathbf{s}'\) belong to D. For any integer \(j\) with \(0 \leqslant j \leqslant q + 1\) define a sequence \(\mathbf{s}_j\) of \(q\) elements of S as follows:
\end{enumerate}

\[
\begin{array}{r l} \mathbf {s} _ {0} & = (s _ {1} ^ {\prime}, \dots , s _ {q} ^ {\prime}) \\ \mathbf {s} _ {1} & = (s _ {1}, \dots , s _ {q}) \\ \mathbf {s} _ {q + 1 - k} & = (s _ {1}, s _ {1} ^ {\prime}, \dots , s _ {1}, s _ {1} ^ {\prime}, s _ {1}, s _ {2}, \dots , s _ {k}) \\ & \text {for } q - k \text { even and } 0 \leqslant k \leqslant q \end{array}\tag{17}
\]

\[
\mathbf {s} _ {q + 1 - k} = \left(s _ {1} ^ {\prime}, s _ {1}, \ldots , s _ {1}, s _ {1} ^ {\prime}, s _ {1}, s _ {2}, \ldots , s _ {k}\right)
\]

\[
\text { for   } q - k \text {   odd   and   } 0 \leqslant k \leqslant q.
\]

Denote by \((\mathrm{H}_{j})\) the assertion `` \(s_{j} \in D, s_{j+1} \in D\) and \(F(s_{j}) \neq F(s_{j+1})\) ''. By A), \((\mathrm{H}_{j}) \Longrightarrow (\mathrm{H}_{j+1})\) for \(0 \leqslant j < q\) , and \((\mathrm{H}_{q})\) is not satisfied by condition b). Hence, \((\mathrm{H}_{0})\) is not satisfied. Since \(s_{0} = s'\) and \(s_{1} = s\) , it follows that \(F(s) = F(s')\) .

PROPOSITION 5. Let M be a monoid (with unit element 1) and f a map from S to M. For \(s, s' \in S\) , let \(m(s, s')\) be the order of \(ss'\) and put

\[
a (s, s ^ {\prime}) = \left\{ \begin{array}{l l} (f (s) f (s ^ {\prime})) ^ {l} & \text {if } m(s,s^{\prime}) = 2l,\ l \text { finite} \\ (f (s) f (s ^ {\prime})) ^ {l} f (s) & \text {if } m(s,s^{\prime}) = 2l + 1,\ l \text { finite} \\ 1 & \text {if } m(s,s^{\prime}) = \infty . \end{array} \right.\tag{18}
\]

If \(a(s, s') = a(s', s)\) whenever \(s \neq s'\) are in \(S\) , there exists a map \(g\) from \(W\) to \(M\) such that

\[
g (w) = f \left(s _ {1}\right) \dots f \left(s _ {q}\right)\tag{19}
\]

for all \(w \in \mathbf{W}\) and any reduced decomposition \((s_1, \ldots, s_q)\) of \(w\) .

For any \(w \in W\) , let \(D_w\) be the set of reduced decompositions of \(w\) and \(F_w\) the map from \(D_w\) to \(M\) defined by

\[
\mathrm{F} _ {w} (s _ {1}, \dots , s _ {q}) = f (s _ {1}) \dots f (s _ {q}).
\]

We are going to prove by induction on the length of w that \(F_{w}\) is constant, which will establish Prop. 5. The cases \(l(w) = 0, 1\) being trivial, we assume that \(q \geqslant 2\) and that our assertion is proved for the elements w with \(l(w) < q\) . Let w be of length q and \(s, s' \in D_w\) ; by Lemma 4 it suffices to prove that \(\mathrm{F}_{w}(\mathbf{s}) = \mathrm{F}_{w}(\mathbf{s}')\) in cases a) and b) of that lemma.

\begin{enumerate}
\def\labelenumi{\alph{enumi})}
\tightlist
\item
  The formula
\end{enumerate}

\[
\mathrm{F} _ {w} \left(s _ {1}, \dots , s _ {q}\right) = f \left(s _ {1}\right) \mathrm{F} _ {w ^ {\prime \prime}} \left(s _ {2}, \dots , s _ {q}\right) = \mathrm{F} _ {w ^ {\prime}} \left(s _ {1}, \dots , s _ {q - 1}\right) f \left(s _ {q}\right)
\]

for \(w' = s_1 \ldots s_{q-1}\) and \(w'' = s_2 \ldots s_q\) and the induction hypothesis show that \(F_w(\mathbf{s}) = F_w(\mathbf{s}')\) if \(s_1 = s_1'\) or \(s_q = s_q'\) .

\begin{enumerate}
\def\labelenumi{\alph{enumi})}
\setcounter{enumi}{1}
\tightlist
\item
  Suppose that there exist two elements \(s\) and \(s'\) of \(S\) such that \(s_j = s'_k = s\) and \(s_k = s'_j = s'\) for \(j\) odd and \(k\) even. It suffices to treat the case \(s \neq s'\) . The sequences \(s\) and \(s'\) are then two distinct reduced decompositions of \(w\) in the dihedral group generated by \(s\) and \(s'\) . By the Remark in no. 2, the order \(m\) of \(ss'\) is necessarily finite and, in the notation of that remark, \(s = s_m\) and \(s' = s'_m\) . Consequently, \(F_w(s) = a(s, s')\) and \(F_w(s') = a(s', s)\) and hence
\end{enumerate}

\[
\mathrm{F} _ {w} (\mathbf {s}) = \mathrm{F} _ {w} (\mathbf {s} ^ {\prime}).
\]

